\section{Background}
\label{sec:background}

This section establishes the technical foundations on which the
methodology builds. We define RAG pipeline components, chunking
strategies, faithfulness measurement, and the Context Coherence Score.

\subsection{Retrieval-Augmented Generation}

RAG systems~\citep{lewis2020rag} augment language model generation with
external documents retrieved at inference time. The standard pipeline
operates in three stages: (1)~\emph{indexing}---documents are split into
chunks and embedded in a vector store; (2)~\emph{retrieval}---the top-$k$
chunks by similarity to the input query are selected; and
(3)~\emph{generation}---the retrieved chunks are concatenated into a prompt,
and the model generates an answer conditioned on both the query and the
retrieved context.

Dense retrieval~\citep{karpukhin2020dpr} replaced lexical matching with
learned embeddings, and cross-encoder re-ranking~\citep{nogueira2019passage,
glass2022re2g} adds a second scoring stage that reorders an over-fetched
candidate set. More recent work explores context
compression~\citep{xu2023recomp} and self-reflective
retrieval~\citep{asai2024selfrag}. All of these interventions target
\emph{retrieval quality} — the relevance of individual chunks to the
query. What they do not directly target, and what this paper studies,
is the \emph{coherence} of the retrieved context as a whole.

\subsection{Chunking and its downstream consequences}

How documents are split determines what the retriever can return.
Fixed-length chunking with a sliding
window~\citep{kamradt2024chunking} is the most common approach, splitting
documents into segments of a fixed token count (e.g., 256, 512, or 1024)
with small overlap. The choice of chunk size creates a fundamental
trade-off: smaller chunks are more precisely matchable to queries but contain
less surrounding context for the generator; larger chunks preserve
intra-document coherence but may include irrelevant material that dilutes
the signal.

Despite its centrality, the faithfulness impact of chunk size has not
been quantified in controlled settings. Most benchmarking work evaluates
chunk size through retrieval accuracy rather than answer
faithfulness~\citep{chen2024benchmarking}. We provide a controlled
measurement of how chunk size affects faithfulness across datasets and
models, and show that it dominates other configuration parameters by a
wide margin.

\subsection{Faithfulness and hallucination}

A generated answer is \emph{faithful} if every claim it makes is supported by
the provided context, regardless of whether those claims are factually
correct~\citep{maynez2020faithfulness}. Hallucination occurs when the model
generates content not grounded in the context, typically by drawing on its
parametric memory~\citep{ji2023hallucination, xu2024hallucination}.

NLI-based faithfulness scoring~\citep{es2024ragas, min2023factscore} treats
the retrieved context as the premise and the generated answer as the
hypothesis. A natural language inference model classifies each answer sentence
as entailed, neutral, or contradicted. The proportion of entailed sentences
yields a faithfulness score in $[0, 1]$. We use
\texttt{cross-encoder/nli-deberta-v3-base}~\citep{he2021deberta} with a
hallucination threshold of 0.50: any response scoring below this threshold
is classified as hallucinated.

\subsection{Context Coherence Score (CCS)}
\label{sec:ccs_def}

We introduce the Context Coherence Score as a retrieval-time diagnostic for
the internal consistency of the retrieved context. The design is motivated by
a specific observation about how generators fail: hallucination tends to
occur not when every passage is poor, but when the retrieved set contains a
mixture of closely related and loosely related passages that the model cannot
reconcile into a single coherent narrative.

\paragraph{Definition.}
Given $n$ retrieved chunks with embeddings $\{e_1, \ldots, e_n\}$, let $S$
denote the matrix of pairwise cosine similarities, where
$S_{ij} = \cos(e_i, e_j)$ for $i \neq j$. We define:
\begin{equation}
  \text{CCS} = \bar{S} - \sigma(S),
  \label{eq:ccs}
\end{equation}
where $\bar{S}$ is the mean of all off-diagonal pairwise similarities and
$\sigma(S)$ is their standard deviation.

\paragraph{Intuition.}
CCS measures whether the retrieved passages tell the same story.
A high CCS means the passages are tightly clustered and uniformly
relevant; a low CCS means the generator is handed a patchwork of
loosely related fragments that it must reconcile on its own.

\paragraph{Design rationale.}
The two terms in CCS capture complementary requirements for a usable
evidence set. The mean $\bar{S}$ measures \emph{topical concentration}: a
high value indicates that all retrieved passages occupy a similar region of
the embedding space, suggesting they address the same subject matter. The
standard deviation $\sigma(S)$ measures \emph{uniformity}: a high value
indicates that some passage pairs are much more similar than others,
signaling an uneven evidence set where strong and weak passages coexist.

Both conditions are necessary. A high mean with high variance describes a
context where one dominant passage is highly similar to the query while
others drift---exactly the setting where generators are most prone to
hallucination, because the model receives conflicting signals about what the
evidence supports. The subtraction $\bar{S} - \sigma(S)$ rewards contexts
that are both topically concentrated \emph{and} internally uniform.

This formulation connects to two established principles. First, in
clustering, cluster quality depends jointly on cohesion (average
intra-cluster similarity) and uniformity of pairwise
distances~\citep{reimers2019sbert}. A tight, uniform cluster of retrieved
passages is analogous to a high-quality cluster whose members agree. A
dispersed, uneven set is analogous to a poorly formed cluster that happens
to contain a few well-placed points. Second, the subtraction of standard
deviation from the mean is structurally analogous to the Sharpe ratio in
portfolio theory, which penalizes high variance in returns even when the
mean return is high. In our setting, high-variance similarity indicates
that the generator receives unevenly grounded evidence, creating
opportunities for hallucination even when the average relevance is adequate.

\paragraph{Scope and limitations.}
CCS is deliberately designed as a lightweight \emph{diagnostic}, not as a
replacement for answer-level faithfulness evaluation. It captures a necessary
condition for faithful generation (coherent evidence) but not a sufficient
one (the answer must also correctly use that evidence). Its value is
operational: CCS can be computed at retrieval time in $O(n^2)$ using only
the embedding model, without running the generator or the NLI scorer. For
typical $k \leq 10$, this adds negligible overhead. The metric serves as a
pre-generation guard-rail---flagging contexts that are likely to cause
problems before committing to expensive generation and evaluation.
